%==============================================================================
%  10-operating-models.tex — Organizational Operating Models
%==============================================================================
\strategyPart{12 \textbar\ Operating Models}{How to Organize for AI}

The operating model decides whether AI scales or stalls. Five archetypes
recur in the research [Deloitte; McKinsey; enterprise-architecture practice].

\section{The five operating models}

\begin{center}
  \renewcommand{\arraystretch}{1.45}
  \rowcolors{2}{inSnow}{white}
  \fitwidth{\begin{tabular}{>{\raggedright\arraybackslash}p{2.4cm}>{\raggedright\arraybackslash}p{4.0cm}>{\raggedright\arraybackslash}p{3.0cm}>{\raggedright\arraybackslash}p{2.6cm}}
    \rowcolor{inNavy}
    \hcell{Model} & \hcell{Main characteristics} & \hcell{Strengths} & \hcell{Weaknesses} \\
    \textbf{Centralized CoE} & One central AI team owns standards, platform, governance & Consistent standards; strong control & Slow, distant from business \\
    \textbf{Federated} & Business units own AI; enterprise principles only & Fast, domain-close & Duplication; inconsistent governance \\
    \textbf{Hybrid} & Central platform + embedded product teams & Balance of control and speed & Needs clear decision rights \\
    \textbf{Product-based} & AI managed as products with owners and roadmaps & User- and value-focused & Requires mature product discipline \\
    \textbf{Agent-orchestrated} & Agents coordinate work and tools under governance & Highest automation potential & Complex governance and security \\
  \end{tabular}}
\end{center}

\section{The AI product operating model}

The product-based model deserves emphasis because it corrects the most common
structural failure: treating AI as a project. Each AI product should have
defined roles:

\begin{itemize}
  \item \textbf{Product owner} — value and roadmap.
  \item \textbf{Business sponsor} — budget and outcomes.
  \item \textbf{Technical owner} — engineering and operations.
  \item \textbf{Data owner} — quality and access.
  \item \textbf{Risk owner} — governance and compliance.
  \item \textbf{User group} — adoption and feedback.
\end{itemize}

\section{Choosing the model}

\begin{itemize}
  \item \textbf{Early maturity} $\rightarrow$ centralized: control and
        standards come first.
  \item \textbf{Diversified enterprises} $\rightarrow$ federated or hybrid:
        domain ownership matters more than standardization.
  \item \textbf{Strategic AI products} $\rightarrow$ product-based: continuity
        beats project completion.
  \item \textbf{Advanced maturity} $\rightarrow$ agent-orchestrated: autonomy
        requires an orchestration layer with human escalation, audit logs and
        permission boundaries.
\end{itemize}

\section{Agent operations team}

An advanced organization adds an \textbf{agent operations capability}
responsible for: agent inventory, tool permissions, prompt and policy
management, performance and cost monitoring, agent evaluation, incident
response and human escalation [Deloitte 2026; Microsoft agent framework].

\section{Organizational structure and agent model}

\begin{center}
  \fitwidth{\begin{tabular}{>{\raggedright\arraybackslash}p{2.9cm}>{\raggedright\arraybackslash}p{3.3cm}>{\raggedright\arraybackslash}p{2.6cm}>{\raggedright\arraybackslash}p{2.9cm}}
    \rowcolor{inNavy}
    \hcell{Structure} & \hcell{Suitable AI model} & \hcell{Main advantage} & \hcell{Main risk} \\
    \textbf{Traditional hierarchy} & Assistants, copilots & Easy adoption & Limited transformation \\
    \textbf{Centralized digital team} & Shared AI services & Consistency & Distance from business \\
    \textbf{Federated units} & Departmental agents & Domain ownership & Duplication \\
    \textbf{Hybrid product model} & Enterprise AI products & Scalable, user-centered & Complex coordination \\
    \textbf{Networked organization} & Multi-agent ecosystem & Fast execution & Governance complexity \\
  \end{tabular}}
\end{center}

\section{Compatibility of agents with the organization}

An agent should fit the organization on seven dimensions [NIST AI RMF;
EU AI Act; ISO/IEC 42001]:
\begin{itemize}
  \item \textbf{Strategy} — connected to cost, customer experience, revenue,
        quality, innovation, risk or resilience objectives.
  \item \textbf{Processes} — a defined trigger, input, reasoning, tool use,
        approval, execution, output, monitoring and escalation.
  \item \textbf{Roles} — a clear relationship with a user, supervisor,
        approver, data owner, product owner, risk owner and administrator.
  \item \textbf{Governance} — registration, risk classification, approvals,
        data-access rules, testing, monitoring, incident response and
        retirement.
  \item \textbf{Technology} — integration with identity, APIs, databases,
        ERP/CRM, BI, collaboration and workflow systems.
  \item \textbf{Culture} — experimentation, training, transparent
        communication, participation, human accountability and safe error
        reporting.
  \item \textbf{Regulation} — classification by data, customer, employee,
        financial, safety and legal impact under the EU AI Act risk tiers.
\end{itemize}

\takeaway{Pick the operating model by maturity and context; agent-orchestrated models deliver the most automation but require the strongest governance.}
